\chapter{7.9 关于推动供应链金融服务实体经济的指导意见(银保监办)} % 2019

\begin{center}
{\small 银保监办发〔2019〕155号} \\
\end{center}

\noindent 各银保监局，各政策性银行、大型银行、股份制银行，邮储银行，外资银行，各保险集团（控股）公司、保险公司、保险资产管理公司，银行业协会、保险业协会：

为深入贯彻党中央、国务院关于推进供应链创新与应用的决策部署，指导银行保险机构规范开展供应链金融业务，推动供应链金融创新，提升金融服务实体经济质效，进一步改善小微企业、民营企业金融服务，现提出以下意见：

\section{总体要求和基本原则}

{ \bf （一）总体要求}

银行保险机构应依托供应链核心企业，基于核心企业与上下游链条企业之间的真实交易，整合物流、信息流、资金流等各类信息，为供应链上下游链条企业提供融资、结算、现金管理等一揽子综合金融服务。

{ \bf （二）基本原则}

银行保险机构在开展供应链金融业务时应坚持以下基本原则：一是坚持精准金融服务，以市场需求为导向，重点支持符合国家产业政策方向、主业集中于实体经济、技术先进、有市场竞争力的产业链链条企业。二是坚持交易背景真实，严防虚假交易、虚构融资、非法获利现象。三是坚持交易信息可得，确保直接获取第一手的原始交易信息和数据。四是坚持全面管控风险，既要关注核心企业的风险变化，也要监测上下游链条企业的风险。

\section{范创新供应链金融业务模式}

{ \bf （三）提供全产业链金融服务}

鼓励银行业金融机构在充分保障客户信息安全的前提下，将金融服务向上游供应前端和下游消费终端延伸，提供覆盖全产业链的金融服务。应根据产业链特点和各交易环节融资需求，量身定制供应链综合金融服务方案。

{ \bf （四）依托核心企业}

鼓励银行业全融机构加强与供应链核心企业的合作，推动核心企业为上下游链条企业增信或向银行提供有效信息，实现全产业链协同健康发展。对于上游企业供应链融资业务，推动核心企业将账款直接付款至专户。对于下游企业供应链融资业务，推动核心企业协助银行整合“三流”信息，并合理承担担保、回购、差额补足等责任。

{ \bf （五）创新发展在线业务}

鼓励银行业金融机构在依法合规、信息交互充分、风险管控有效的基础上，运用互联网、物联网、区块链、生物识别、人工智能等技术，与核心企业等合作搭建服务上下游链条企业的供应链金融服务平台，完善风控技术和模型，创新发展在线金融产品和服务，实施在线审批和放款，更好满足企业融资需求。

{ \bf （六）优化结算业务}

银行业金融机构应根据供应链上下游链条企业的行业结算特点，以及不同交易环节的结算需求，拓展符合企业实际的支付结算和现金管理服务，提升供应链支付结算效率。

{ \bf （七）发展保险业务}

保险机构应根据供应链发展特点，在供应链融资业务中稳妥开展各类信用保证保险业务，为上下游链条企业获取融资提供增信支持。

{ \bf （八）加强小微民营企业金融服务}

鼓励银行保险机构加强对供应链上下游小微企业、民营企业的金融支持，提高金融服务的覆盖面、可得性和便利性，合理确定贷款期限，努力降低企业融资成本。

{ \bf （九）加强“三农”金融服务}

鼓励银行保险机构开展农业供应链金融服务和创新，支持订单农户参加农业保险，将金融服务延伸至种植户、养殖户等终端农户，以核心企业带动农村企业和农户发展，促进乡村振兴。

\section{完善供应链金融业务管理体系}

{ \bf （十）加强业务集中管理}

鼓励银行保险机构成立供应链金融业务管理部门（中心），加强供应链金融业务的集中统一管理，统筹推进供应链金融业务创新发展，加快培育专业人才队伍。

{ \bf （十一）合理配置供应链融资额度}

银行业金融机构应合理核定供应链核心企业、上下游链条企业的授信额度，基于供应链上下游交易情况，对不同主体分别实施额度管理，满足供应链有效融资需求。其中，对于由核心企业承担最终偿付责任的供应链融资业务，应全额纳入核心企业授信进行统一管理，并遵守大额风险暴露的相关监管要求。

{ \bf （十二）实施差别化信贷管理}

在有效控制风险的前提下，银行业金融机构可根据在线供应链金融业务的特点，制定有针对性的信贷管理办法，通过在线审核交易单据确保交易真实性，通过与供应链核心企业、全国和地方信用信息共享平台等机构的信息共享，依托工商、税务、司法、征信等数据，采取在线信息分析与线下抽查相结合的方式，开展贷款“三查”工作。

{ \bf （十三）完善激励约束机制}

银行保险机构应健全供应链金融业务激励约束及容错纠错机制，科学设置考核指标体系。对于供应链上下游小微企业贷款，应落实好不良贷款容忍度、尽职免责等政策。

{ \bf （十四）推动银保合作}

支持银行业金融机构和保险机构加强沟通协商，在客户拓展、系统开发、信息共享、业务培训、欠款追偿等多个环节开展合作。协同加强全面风险管理，共同防范骗贷骗赔风险。

\section{加强供应链金融风险管控}

{ \bf （十五）加强总体风险管控}

银行业金融机构应建立健全面向供应链金融全链条的风险控制体系，根据供应链金融业务特点，提高事前、事中、事后各个环节的风险管理针对性和有效性，确保资金流向实体经济。

{ \bf （十六）加强核心企业风险管理}

银行业金融机构应加强对核心企业经营状况、核心企业与上下游链条企业交易情况的监控，分析供应链历史交易记录，加强对物流、信息流、资金流和第三方数据等信息的跟踪管理。银行保险机构应明确核心企业准入标准和名单动态管理机制，加强对核心企业所处行业发展前景的研判，及时开展风险预警、核查与处置。

{ \bf （十七）加强真实性审查}

银行业金融机构在开展供应链融资业务时，应对交易真实性和合理性进行尽职审核与专业判断。鼓励银行保险机构将物联网、区块链等新技术嵌入交易环节，运用移动感知视频、电子围栏、卫星定位、无线射频识别等技术，对物流及库存商品实施远程监测，提升智能风控水平。

{ \bf （十八）加强合规管理}

银行保险机构应加强供应链金融业务的合规管理，切实按照回归本源、专注主业的要求，合规审慎开展业务创新，禁止借金融创新之名违法违规展业或变相开办未经许可的业务。不得借供应链金融之名搭建提供撮合和报价等中介服务的多边资产交易平台。

{ \bf （十九）加强信息科技系统建设}

银行保险机构应加强信息科技系统建设，鼓励开发供应链金融专项信息科技系统，加强运维管理，保障数据安全，借助系统提升风控技术和能力。

\section{优化供应链金融发展的外部环境}

{ \bf （二十）加强产品推介}

银行保险机构应加强供应链金融产品的开发与推介，及时宣传供应链金融服务小微企业、民营企业进展情况。

{ \bf （二十一）促进行业交流}

银行业和保险业自律组织应组织推动行业交流，总结推广银行保险机构在供应链金融领域的良好实践和经验，促进供应链金融持续健康发展。

{ \bf （二十二）提高监管有效性}

各级监管部门应根据供应链金融业务特点，加强供应链金融风险监管，规范银行业金融机构和保险机构的业务合作，对于业务经营中的不审慎和违法违规行为，及时采取监管措施。 \\

\rightline{2019年7月9日}
